\chapter{Az admin és a publikus kliens megvalósítása}

\section{Admin (React, TypeScript, Vite)}
Az admin SPA moduljai: bejelentkezés, felhasználók kezelése (csak \texttt{admin} szerepkörrel elérhető menü), szólisták és színpaletták. A token és a felhasználó \texttt{localStorage}-ban tárolódik (\texttt{cellauto\_admin\_token}, \texttt{cellauto\_admin\_user}). A hibák kinyerése több JSON mezőből történik (\texttt{error}, \texttt{message}, Laravel \texttt{errors}).

\section{A \texttt{www} kliens}
Vanilla JS modulok: \texttt{script.js} (rács, szabályok, generációk), \texttt{app.js} (bejelentkezés, lista- és színbetöltés, dinamikus CSS a hex színekhez), \texttt{board-save.js} (mentési csoportok és slotok UI), \texttt{api.js} (fetch wrapper). A PHP oldal routere a statikus és API-proxy útvonalakat szolgálja ki a környezettől függően.

\chapter{Telepítés, tesztelés, biztonság}

\section{Telepítés}
\texttt{composer run setup} az API repóban: függőségek, kulcs, migráció, npm build. Fejlesztői mód: \texttt{composer run dev}. Éles: \texttt{.env}, \texttt{php artisan migrate --force}, opcionálisan queue worker.

\section{Tesztelés}
\texttt{composer run test} — PHPUnit tesztek. Az \texttt{implementacios-terv.md} további feature teszteket javasol (auth, board CRUD folyamat).

\section{Biztonság}
Sanctum tokenek; admin middleware; HTTPS; élesben CSP és XSS védelem; jelszavak hash-elve (\texttt{bcrypt}).

\chapter{Összegzés és továbbfejlesztés}

\section{Eredmény}
A dolgozat egy összhangban lévő háromrétegű rendszert mutat be elméleti indokkal, részletes adatmodellel és teljes REST áttekintéssel.

\section{Irányok}
Form Request validáció, egységes hiba JSON, lapozás a listáknál, board lista válasz payload nélkül opcionális query-vel, reorder végpontok, bővített tesztlefedettség.

\begin{thebibliography}{9}
\bibitem{wolfram2002}
\textsc{Wolfram, S.}: \emph{A New Kind of Science}, Wolfram Media, 2002.
\bibitem{laravel}
\textsc{Laravel dokumentáció}, \url{https://laravel.com/docs}
\bibitem{react}
\textsc{React dokumentáció}, \url{https://react.dev}
\bibitem{cellauto-docs}
\textsc{Cellauto projekt dokumentáció}, \texttt{admin/docs}, \texttt{api/docs}, \texttt{www/docs} könyvtárak.
\end{thebibliography}

\appendix
\chapter{Forráskód és dokumentáció helye}
A \texttt{collect-sources.sh} szkript a három alkalmazás másolatát egy mappába gyűjti (\texttt{UTMUTATO.txt}). A végpontok és sémák változásakor a markdown dokumentációt célszerű a kóddal együtt frissíteni.
